\chapter{The Six Knobs That Run the Sky}

\epi{Credit where it's due: six dials, thousands of numbers, one of the best fits in the history of science.\\ Then the fine print.}

\fig[1.0]{plate_cmb_sky}{The sky the six knobs describe. A simulated patch of microwave sky made from a toy power spectrum (not real data): hot and cold spots that differ by about a hundred-thousandth of the average temperature. Underneath, the six dials that set its statistics.}


\section{A universe on six dials}

If you want to see knob economy at its finest, look at the cosmic microwave background. This is the afterglow of the hot early universe, released about 380,000 years after the beginning when the cosmos cooled enough for electrons and protons to combine into neutral hydrogen and let light travel freely. The Planck satellite mapped its temperature and polarisation across the whole sky. Statistically, the map is summarised by power spectra: how much the temperature varies, on average, on each angular scale. The spectra are labelled by multipole $\ell$, where $\ell\approx180^\circ/\theta$ for an angle $\theta$. Planck measured them from $\ell=2$ up to $\ell\approx2500$, in temperature (TT), polarisation (EE) and their cross-correlation (TE). That's thousands of measured numbers.

The spectrum has a shape you can learn to read, like a fingerprint. Each feature is set mainly by one or two of the knobs below, which is why six numbers can pin down thousands of measurements.


\fig{cmb_anatomy}{Reading the microwave spectrum (a cartoon). The large scales carry the primordial seeds; the first peak's position is the acoustic angle; the second peak is lowered by baryons loading the sound waves; the third is driven by dark matter's gravity; and the tail is erased by photons diffusing.}



Base \LCDM\ fits all of them with six knobs. I'll introduce them properly, because in a book about knobs these six deserve a bow before the jokes start.

\section{The expansion: Friedmann's equation}

Everything rests on one equation from general relativity applied to a uniform universe. The expansion rate $H$ at redshift $z$ obeys
\begin{equation}
H^2(z)=H_0^2\left[\Omega_r(1+z)^4+\Omega_m(1+z)^3+\Omega_k(1+z)^2+\Omega_\Lambda\right],
\end{equation}
where $H_0$ is today's expansion rate and each $\Omega$ is a component's share of the critical density $\rho_c=3H_0^2/8\pi G$: radiation (falling as $(1+z)^4$), matter ($(1+z)^3$), spatial curvature ($(1+z)^2$) and the cosmological constant (constant). In base \LCDM\ the universe is flat, $\Omega_k=0$, so $\Omega_\Lambda=1-\Omega_m-\Omega_r$ is not a separate knob. People sometimes say \LCDM\ has seven parameters by counting $\Lambda$ twice. It has six.

\fig{friedmann}{Friedmann's equation, one term at a time. Radiation dominates the early universe, matter takes over at $z\approx3400$, and the cosmological constant takes over only at $z\approx0.3$. The CMB was released in the matter era. (Densities from the standard reading of the expansion history, for illustration.)}


\section{The six, one at a time}

\begin{description}[leftmargin=1em,style=nextline]
\item[1. $\omega_b=\Omega_bh^2\approx0.0224$, the baryon density.] Ordinary matter (protons, neutrons, the stuff of stars and you), with $h=H_0/(100\kms\,\mathrm{Mpc}^{-1})$. Baryons add inertia to the photon fluid. The ratio $R=3\rho_b/4\rho_\gamma$ shifts the balance between gravity and pressure, which makes the compressions (odd-numbered peaks) higher than the rarefactions (even-numbered peaks). The ratio of the first two peaks reads off $\omega_b$ almost directly. It agrees with the deuterium abundance from Big Bang nucleosynthesis, an independent physical process minutes after the beginning. That agreement is one of the great cross-checks in physics.
\item[2. $\omega_c=\Omega_ch^2\approx0.120$, the cold dark matter density.] This is about five times the baryon density. It sets when matter came to outweigh radiation, which controls how strongly the oscillations are driven, and it shapes the third peak, which is sensitive to potential wells that don't oscillate with the photons. It is the amount of a substance that has never been detected in a laboratory.
\item[3. $\theta_*\approx1.041\times10^{-2}$, the acoustic angle.] The sound horizon at last scattering divided by the distance to it:
\begin{equation}
\theta_*=\frac{r_s(z_*)}{D_M(z_*)},\qquad r_s=\int_{z_*}^{\infty}\frac{c_s(z)\,dz}{H(z)},\qquad c_s=\frac{c}{\sqrt{3(1+R)}}.
\end{equation}
It fixes where the peaks fall, and it's measured to about 0.03 per cent, one of the most precise numbers in cosmology. Combined with the others it gives $H_0\approx67.4\kms\,\mathrm{Mpc}^{-1}$.
\item[4. $\tau\approx0.054$, the optical depth to reionisation.] When the first stars and quasars reionised the universe, free electrons re-scattered some CMB light. That suppresses the peaks by a factor of $e^{-2\tau}$ on small scales and creates a bump in polarisation on the largest scales.
\item[5 and 6. $A_s\approx2.1\times10^{-9}$ and $n_s\approx0.965$, the primordial fluctuations.] The initial seeds are assumed to follow a power law,
\begin{equation}
\mathcal{P}_\mathcal{R}(k)=A_s\left(\frac{k}{k_{\rm p}}\right)^{n_s-1},\qquad k_{\rm p}=0.05\,\mathrm{Mpc}^{-1},
\end{equation}
with amplitude $A_s$ and tilt $n_s$. That $n_s$ comes out slightly below 1 is a genuine success for inflation, which predicted a small tilt.
\end{description}

\fig{cmb_cartoon}{What four of the six knobs do to the CMB temperature spectrum. \textbf{This is a cartoon}, a toy oscillation with damping, not a computed spectrum. Baryons raise the odd peaks relative to the even ones; the acoustic angle slides the peaks; reionisation lowers everything above the largest scales; the tilt tips small scales up or down.}


\begin{knobbox}{Base \LCDM\ at the CMB}
Six global knobs: $\omega_b$, $\omega_c$, $\theta_*$, $\tau$, $A_s$, $n_s$. Thousands of measured spectrum values fitted, and the predicted shape of each spectrum carries information on all six at once. By the information measure of Chapter~1, each knob is pinned at the per-cent level or better, which is roughly 6 to 12 bits per knob, on the order of 50 bits for the set. That's an astonishingly good exchange rate.
\end{knobbox}

The acoustic angle deserves a picture of its own, because it is the most precise number in cosmology and it works like a ruler held at arm's length.


\fig{sound_horizon}{The sound horizon as a ruler. Left: a sound wave spreads from every clump in the early universe until the light is released; the distance it has travelled then is $r_s$. Right: seen from across the observable universe, that ruler spans the angle $\theta_*=r_s/D_M$, about 0.6 degrees on the sky.}



\begin{honestbox}[Honesty box: this is the theory's strongest ground]
The CMB is exactly where my own law has nothing to offer. The acceleration scale $a_0$ doesn't appear in the linear equations of a uniform expanding universe, and the programme's record has no CMB fit without a component that behaves, at linear order, like cold dark matter. When the law was run alone on the cosmic web, starting from baryons only, it got the shape of the CMB lensing spectrum wrong under every rule tried (Chapter~12). Any reader who comes away from this book thinking that knob-counting has overturned \LCDM's cosmology has misread it.
\end{honestbox}

\section{The fine print: knobs behind the knobs}

Now the jokes can start. Six is the number on the box. Open it and there are more.

\subsection{Nuisance knobs}

The microwave sky isn't just the CMB. There's dust in our galaxy, radio sources, the cosmic infrared background from distant star-forming galaxies, and the Sunyaev--Zel'dovich effect from hot gas in clusters. Each detector has calibration and beam uncertainties. A real CMB likelihood carries its own nuisance knobs for all of this, of order twenty or more in the high-multipole temperature-and-polarisation likelihood. Most are constrained by priors or by other frequency channels, and they're marginalised over, so they're not freedom in the worrying sense. They are still knobs, and they're the reason ``six parameters'' should always be read as ``six parameters of physical interest.''

\fig{lcdm_panel_full}{The full \LCDM\ control panel. Top: the six. Middle: nine dials left at default settings, any of which can be freed. Bottom: the nuisance dials for foregrounds and calibration. Only the top row is usually quoted.}



\subsection{Hidden knobs: the choices}

The six are six \emph{given} a set of choices: flat geometry, three neutrino species with a minimal mass sum, a helium abundance from nucleosynthesis, a power-law spectrum with no running, no primordial gravitational waves, dark energy as a pure constant. Each of those is a knob fixed at a default value. Freeing them gives the well-known extensions:
\begin{equation}
\{\Omega_k,\ \textstyle\sum m_\nu,\ N_{\rm eff},\ Y_P,\ dn_s/d\ln k,\ r,\ w_0,\ w_a,\ A_{\rm lens},\dots\}
\end{equation}
Each extension is a dial that \LCDM\ can turn if the data demand it. That's legitimate; it's what a framework with a good base model should do. It does mean that the honest count of what the theory \emph{could} accommodate is larger than six. In the language of the Occam factor, the prior is wider than the base model suggests.

One of these deserves a wink. $A_{\rm lens}$ scales the amount of gravitational-lensing smoothing of the CMB peaks. In a consistent theory it must be 1, so it isn't physics but a consistency check. In the 2018 Planck analysis the temperature data preferred a value somewhat above 1. When a knob that is supposed to be fixed at 1 wants to be turned, that's the universe tapping on the glass.

\subsection{The knob with no theory}

$\Omega_\Lambda$ is fixed by flatness, but $\Lambda$ itself is the most embarrassing knob in physics. If it's the energy of the vacuum, a naive quantum estimate with a cutoff at the Planck scale overshoots the measured value by a factor of about $10^{120}$. Even with a much lower cutoff, the mismatch is dozens of orders of magnitude. Nothing in \LCDM\ explains the value; it's simply measured. And it shows up just when it's comparable to the matter density, which is the ``why now?'' coincidence. Keep an eye on this one: my galaxy law ties its acceleration scale to exactly this density, so it inherits the embarrassment rather than curing it.

How big is the embarrassment? It depends where you stop counting the vacuum's energy, but every choice is absurd. Stopping at the scale of the strong nuclear force gives a vacuum energy about $10^{44}$ times too big; stopping at the Planck scale, about $10^{123}$ times.


\fig{vacuum_catastrophe}{The worst estimate in physics. A naive vacuum energy, cut off at four energy scales, against the measured dark energy density. Even the most modest cutoff overshoots by 44 powers of ten.}



\section{Where the six start to creak}

A knob set in one place should give the same value when measured somewhere else. When it doesn't, either the measurement is wrong or the knob is standing in for missing physics. Three such tensions are worth knowing.

\textbf{The Hubble tension.} The CMB, through $\theta_*$ and the model, gives $H_0\approx67$--$68\kms\,\mathrm{Mpc}^{-1}$. The local distance ladder of Cepheids and type Ia supernovae gives about $73$, with errors near one unit. That's a gap of about five standard deviations, persistent for a decade. Either some ladder systematic remains, or one of the default knobs (the early expansion history, perhaps) needs turning.

\fig{h0_tension}{The Hubble tension: the same knob read two ways. The CMB with base \LCDM\ gives about 67.4; the local distance ladder gives about 73.}


\textbf{The clumpiness tension.} Weak lensing surveys tend to find the matter distribution slightly less clumpy (a lower $S_8$) than the CMB-calibrated model predicts. Its significance has varied from survey to survey and has eased in some recent analyses.

\textbf{Dark energy that moves.} In 2024 and 2025 the DESI survey's baryon acoustic oscillation measurements, combined with the CMB and supernovae, preferred dark energy whose density changes with time, parameterised by $w_0>-1$ and $w_a<0$ in
\begin{equation}
w(a)=w_0+w_a(1-a),
\end{equation}
over a pure constant, at roughly 3 to 4 standard deviations depending on the supernova sample. If that holds up, the knob called $\Lambda$ was never a constant, and two more dials ($w_0$, $w_a$) move from ``optional'' to ``required.''

\fig{w0wa_plane}{The dark energy knob plane. The red dot is a cosmological constant, where the law's $a_0$ is exactly constant in time. The shaded quadrant is where recent combined fits have leaned (schematic, no contours drawn). The gold square is the illustrative case worked out in Chapter~13.\label{fig:w0wa}}



That last one matters enormously for this book. In my law the galaxy acceleration scale is tied to the dark energy density, so if dark energy moves, the law says galaxies must notice. We'll come back to this in Part~V. It's the knob that can't be turned.

\section{The scorecard so far}

\begin{center}
\begin{tabular}{lrl}
\toprule
item & knobs & comment\\
\midrule
base physics & 6 & superb fit to thousands of CMB numbers\\
nuisance & $\sim$20+ & mostly prior-constrained; not physics\\
default choices & $\sim$9 & each a dial left at a default\\
$\Lambda$'s value & 1 & measured, unexplained by $10^{120}$\\
dark matter's identity & -- & not a number; a missing particle\\
\bottomrule
\end{tabular}
\end{center}

On the cosmological scale, that's a lean and well-run panel. The trouble starts when you zoom in by a factor of a thousand, from the horizon to a galaxy. There, the six knobs say nothing about what a galaxy should look like, and a whole new panel has to be bolted on. That's next.
